\scanpage{169}

UNIVERSITE DES SCIENCES

ET TECHNIQUES DU LANGUEDOC

M A T H E M A T I Q U E S \hfill 21.10.80

\begin{center}
SEMINAIRE D'ALGEBRE

1980-1981
\end{center}

Vendredi de 10 h à 12 h, Salle 431 du Bâtiment de Mathématiques.

31 Octobre, J. HELMSTETTER (Institut Fourier, Université de Grenoble I)

\uline{Algèbres de Clifford et groupes de Clifford.}

\uline{Une sous-variété remarquable $XE$ dans l'algèbre extérieure $\unsure{A|\Lambda}E$.}

Soit $E$ un espace vectoriel de dimension finie sur le corps commutatif $K$ ; $XE$ est l'ensemble des éléments de $\Lambda E$ de la forme $x\wedge\operatorname{Exp}u$, où $x$ est un élément décom-\newline
posable de $\Lambda E$ (il est possible que $x\in K$) et où $u\in\Lambda^2E$ ; le symbole $\operatorname{Exp}$ désigne une exponentielle dans la sous-algèbre paire de $\Lambda E$. On commence par le théorème suivant : le produit intérieur d'un élément de $XE$ par un élément de $XE^*$ (où $E^*$ est le dual de $E$), est encore dans $XE$.

Application : soit $Q$ une forme quadratique sur $E$ et soit $B$ une forme bilinéaire (pas nécessairement symétrique) telle que $Q(a)=B(a,a)$ pour tout $a\in E$ ; on peut ``déformer'' le produit extérieur de l'algèbre $\Lambda E$ pour en faire un nouveau produit (noté $xy$) isomorphe à celui de l'algèbre de Clifford associée à $Q$ ; ainsi l'espace vectoriel $\Lambda E$ devient une algèbre $\Lambda(E,B)$ isomorphe au quotient de l'algèbre tenso-\newline
rielle $TE$ par l'idéal bilatère engendré par les éléments $a\unsure{\otimes}\unsure{a-}Q(a)$ ; si $x$ et $y\in XE$, on démontre que $xy\in XE$.

\uline{Indications sur les exposés suivants} : On étudie l'algèbre $\Lambda(E,B)$ mentionnée ci-dessus

Il existe une application involutive $x\to x'$ de $XE$ dans lui-même telle que $xx'=x'x\in K$ et $xax'\in E$ quels que soient $x\in XE$ et $a\in E$. Le groupe de Clifford $X(E,B)$ est l'ensemble des éléments de $XE$ inversibles dans $\Lambda(E,B)$ ; on définit un homomorphisme surjectif de $X(E,B)$ sur le groupe orthogonal $GO(E,Q)$ (groupe des automorphismes de $E$ qui conservent $Q$ ainsi que tous les éléments de $E$ orthogonaux à $E$).

\hfill A. MICALI

\uline{N.B.} Le Séminaire d'Algèbre se fera, cette année, à une cadence d'un exposé par mois sur \uline{les Algèbres de Clifford et Groupes de Clifford} et de trois à quatre exposés par mois sur les \uline{Automorphismes des anneaux de polynômes et séries formelles}.

\margin{169}
